\documentclass[11pt,a4paper]{article}
\usepackage[utf8]{inputenc}
\usepackage[vietnamese]{babel}
\usepackage{amsmath,amssymb}
\usepackage{geometry}
\geometry{a4paper, total={180mm,260mm}, left=15mm, right=15mm, top=18mm, bottom=18mm}

\usepackage{tikz}
\usetikzlibrary{shapes.geometric, arrows.meta, positioning, calc}
\tikzset{arr/.style={-{Stealth[scale=0.8]}}}
\usepackage{tcolorbox}
\tcbuselibrary{skins,breakable}
\usepackage{fontawesome}
\usepackage{fancyhdr}
\usepackage{enumitem}

\pagestyle{fancy}
\fancyhf{}
\renewcommand{\headrulewidth}{0.5pt}
\renewcommand{\footrulewidth}{0.5pt}
\fancyhead[L]{\small\bfseries\color{violet!85!black}\faBook\ SỔ TAY KIẾN THỨC TOÁN 2}
\fancyhead[R]{\small\bfseries\color{violet!85!black}Lớp Toán Thầy Lĩnh -- loptoanthaylinh.vn}
\fancyfoot[C]{\small\color{gray}-- Sổ tay kiến thức: Bài 02. Tia số. Số liền trước, số liền sau --}

\setlength{\parindent}{0pt}
\setlength{\parskip}{0pt}

\begin{document}

% --- BANNER BÀI HỌC ---
\begin{tcolorbox}[
  enhanced,
  colback=violet!8!white,
  colframe=violet!70!black,
  arc=3mm, boxrule=1.2pt,
  left=3mm, right=3mm, top=2mm, bottom=2mm,
  halign=center
]
  {\LARGE\bfseries\color{violet!90!black}\faBookmark\ BÀI 02: TIA SỐ. SỐ LIỀN TRƯỚC, SỐ LIỀN SAU}\\[1mm]
  {\itshape\footnotesize\color{black!75} Khái niệm tia số, xác định vị trí và quy tắc tìm số liền trước, số liền sau trên 1 trang A4}
\end{tcolorbox}

\vspace{-2mm}

\noindent
\begin{minipage}[t]{0.49\textwidth}\vspace{0pt}
\begin{tcolorbox}[
  enhanced,
  colback=violet!4!white,
  colframe=violet!75!black,
  arc=2mm, boxrule=0.8pt,
  left=2mm, right=2mm, top=1.8mm, bottom=1.8mm,
  title={\footnotesize\bfseries\color{white}\faKey\ 1. Khái niệm tia số}
]
{\footnotesize
\begin{itemize}\setlength{\itemsep}{1pt}\setlength{\leftskip}{-2.5mm}
  \item \textbf{Tia số} là một tia thẳng có gốc là điểm \textbf{0}.
  \item Chiều mũi tên hướng từ \textbf{trái sang phải}, biểu thị các số tăng dần: $0, 1, 2, 3, \dots$
  \item Mỗi vạch trên tia số ứng với một số tự nhiên. Khoảng cách giữa hai vạch liên tiếp luôn \textbf{bằng nhau}.
  \item \textbf{Quy tắc vị trí trên tia số:}\\
  Số ở bên trái luôn \textbf{bé hơn} số ở bên phải.\\
  \textit{Ví dụ:} Số $15$ đứng trước số $18$ nên $15 < 18$.
\end{itemize}
}
\end{tcolorbox}
\end{minipage}
\hfill
\begin{minipage}[t]{0.49\textwidth}\vspace{0pt}
\begin{tcolorbox}[
  enhanced,
  colback=teal!4!white,
  colframe=teal!75!black,
  arc=2mm, boxrule=0.8pt,
  left=2mm, right=2mm, top=1.8mm, bottom=1.8mm,
  title={\footnotesize\bfseries\color{white}\faSliders\ 2. Số liền trước, số liền sau}
]
{\footnotesize
\begin{itemize}\setlength{\itemsep}{1pt}\setlength{\leftskip}{-2.5mm}
  \item \textbf{Số liền trước} của một số là số đứng ngay trước số đó trên tia số (bên trái).
  $$\text{Số liền trước} = \text{Số đã cho} - 1$$
  \item \textbf{Số liền sau} của một số là số đứng ngay sau số đó trên tia số (bên phải).
  $$\text{Số liền sau} = \text{Số đã cho} + 1$$
  \item \textbf{Số ở giữa:} Nếu ba số liên tiếp là $A, B, C$ thì $B$ là số ở giữa hai số $A$ và $C$ ($B = A + 1 = C - 1$).
\end{itemize}
}
\end{tcolorbox}
\end{minipage}

\vspace{-1.5mm}
\begin{tcolorbox}[
  enhanced,
  colback=orange!4!white,
  colframe=orange!80!black,
  arc=2mm, boxrule=0.8pt,
  left=2mm, right=2mm, top=1.5mm, bottom=1.5mm,
  title={\footnotesize\bfseries\color{white}\faSitemap\ 3. Mô hình trực quan tia số và quy tắc cộng/trừ 1}
]
\centering
\begin{tikzpicture}[x=1.2cm, y=1cm]
  % Vẽ tia số
  \draw[->, very thick, violet!80!black] (0,0) -- (6.5,0);
  \foreach \x/\label in {0/18, 1/19, 2/20, 3/21, 4/22, 5/23} {
    \draw[thick, violet!70!black] (\x, 0.15) -- (\x, -0.15) node[below, font=\scriptsize\bfseries] {\label};
  }
  % Nút số 20
  \filldraw[violet!80!black] (2,0) circle (2.5pt) node[above=4pt, font=\scriptsize\bfseries\color{violet!90!black}] {Số 20};
  % Cung liền trước
  \draw[arr, teal!80!black, thick, ->] (2,0.3) to[bend right=40] node[above, font=\tiny\bfseries] {Liền trước: Bớt 1 (-1)} (1,0.3);
  % Cung liền sau
  \draw[arr, orange!80!black, thick, ->] (2,0.3) to[bend left=40] node[above, font=\tiny\bfseries] {Liền sau: Thêm 1 (+1)} (3,0.3);
\end{tikzpicture}
\end{tcolorbox}

\vspace{-1.5mm}
\noindent
\begin{minipage}[t]{0.49\textwidth}\vspace{0pt}
\begin{tcolorbox}[
  enhanced,
  colback=green!4!white,
  colframe=green!60!black,
  arc=2mm, boxrule=0.8pt,
  left=2mm, right=2mm, top=1.8mm, bottom=1.8mm,
  title={\footnotesize\bfseries\color{white}\faCheckCircle\ 4. Ví dụ thực hành chuẩn mực}
]
{\footnotesize
\textbf{1. Tìm số liền trước, liền sau:}\\
\begin{itemize}\setlength{\itemsep}{0.5pt}\setlength{\leftskip}{-2.5mm}
  \item Số liền trước của $40$ là: $40 - 1 = \mathbf{39}$.
  \item Số liền sau của $79$ là: $79 + 1 = \mathbf{80}$.
  \item Số liền trước của $100$ là: $100 - 1 = \mathbf{99}$.
\end{itemize}
\textbf{2. Tìm số ở giữa:}\\
Số ở giữa $59$ và $61$ là số $\mathbf{60}$.
Ba số liên tiếp: $59, 60, 61$.
}
\end{tcolorbox}
\end{minipage}
\hfill
\begin{minipage}[t]{0.49\textwidth}\vspace{0pt}
\begin{tcolorbox}[
  enhanced,
  colback=red!4!white,
  colframe=red!75!black,
  arc=2mm, boxrule=0.8pt,
  left=2mm, right=2mm, top=1.8mm, bottom=1.8mm,
  title={\footnotesize\bfseries\color{white}\faExclamationTriangle\ 5. Mẹo tránh bẫy trắc nghiệm}
]
{\footnotesize
\begin{itemize}\setlength{\itemsep}{1pt}\setlength{\leftskip}{-2.5mm}
  \item \textbf{Bẫy câu chữ: "Liền trước" vs "Liền sau":}\\
  - \textbf{Trước} $\to$ Đứng trước $\to$ Bé hơn $\to$ \textbf{Trừ 1}.\\
  - \textbf{Sau} $\to$ Đứng sau $\to$ Lớn hơn $\to$ \textbf{Cộng 1}.
  \item \textbf{Bẫy số bé nhất có 2 chữ số:}\\
  Số liền trước của số bé nhất có hai chữ số ($10$) là \textbf{9} (số có 1 chữ số).
  \item \textbf{Bẫy mốc tròn chục:}\\
  Liền trước của $90$ là $89$, không phải $91$.
\end{itemize}
}
\end{tcolorbox}
\end{minipage}


\end{document}
